\section{Tóm tắt nghiên cứu}
% [NHÁP -- Phụ trách: Nguyễn Văn Thương. Nhóm rà soát & xác nhận trước khi nộp.]

Chẩn đoán bệnh dựa trên triệu chứng là một bài toán cốt lõi trong lĩnh vực \textbf{Healthcare AI}, nơi việc truy xuất đúng tri thức y khoa có vai trò quyết định đến độ tin cậy và an toàn của hệ thống hỗ trợ ra quyết định lâm sàng. Nghiên cứu này xây dựng và đánh giá một hệ thống \textbf{Retrieval-Augmented Generation (RAG)} hỗ trợ chẩn đoán bệnh, sử dụng bộ dữ liệu công khai \textbf{DDXPlus}~\cite{ddxplus2022} gồm \textbf{49 bệnh} và \textbf{223 bằng chứng triệu chứng (evidences)}. Nhóm làm việc trên \textbf{DDXPlus (English) phiên bản 2, đăng ngày 10/03/2026} trên Figshare.

Chúng tôi xây dựng một Knowledge Base (KB) mô tả bệnh từ DDXPlus, chuẩn hoá và kiểm định chất lượng KB, sau đó so sánh ba chiến lược truy xuất: \textbf{BM25} (lexical), \textbf{Dense} (mô hình \texttt{BAAI/bge-small-en-v1.5}, 384 chiều, độ đo cosine) và \textbf{Hybrid} (hợp nhất thứ hạng bằng Reciprocal Rank Fusion, RRF $k=60$).

\textbf{Lưu ý về phương pháp đo:} các con số EXP-01/EXP-02 trong báo cáo này là kết quả của một \textbf{harness Python thuần, brute-force in-memory} trên 49 tài liệu KB, chạy độc lập với hạ tầng tìm kiếm. Đây là \textit{mốc tham chiếu / trần trên} cho chất lượng logic truy xuất trên cùng KB và normalizer. Số liệu \textit{production} được đo trên \textbf{pipeline OpenSearch thật}: lần chạy đầy đủ ngày \textbf{26/06/2026} trên \textbf{toàn bộ $n=132{,}448$ ca} (thay cho lần chạy thử trên mẫu trước đó) cho kết quả nhất quán với offline -- riêng BM25 tái lập gần như khít (Hit@1 98.65\% production vs 98.47\% offline) -- xác nhận độ tin cậy của cả harness lẫn hạ tầng (xem mục Kết quả và Thảo luận).

Kết quả chính theo Hit@1 (Python thuần, $n=132{,}448$): BM25 + \textit{description} đạt \textbf{98.78\%} (MRR 0.9937); Hybrid-RRF + \textit{description} đạt \textbf{91.95\%} (MRR 0.9570); Dense + \textit{description} đạt \textbf{84.28\%} (MRR 0.9027). Hai phát hiện đáng chú ý: (i) trường \textit{description} cải thiện đều cả ba phương pháp; (ii) thiết kế đánh giá hiện tại thiên về \textit{lexical} (truy vấn là token đã chuẩn hoá, trùng khớp từ vựng với KB), nên lợi thế ngữ nghĩa của dense/hybrid chưa được phản ánh đầy đủ.

\textbf{Đóng góp chính:} (1) xây dựng và kiểm định KB DDXPlus (chuẩn hoá evidence, rà soát ánh xạ mã ICD-10); (2) thiết lập khung đánh giá truy xuất theo Hit@1/3/5, MRR và NDCG@5, kèm quy trình đối chiếu \textbf{Python thuần với OpenSearch trên toàn tập}; (3) so sánh hệ thống truy xuất lexical/dense/hybrid trên quy mô đầy đủ và phân tích trung thực giới hạn của bài đánh giá.